% ============================================================
%  Mangia Consapevole — Guida pratica agli additivi alimentari
%  faustobe.it | 2026 | CC BY-NC 4.0
% ============================================================
\documentclass[a5paper,12pt,oneside,openany]{book}

\usepackage[T1]{fontenc}
\usepackage[utf8]{inputenc}
\usepackage{ebgaramond}
\usepackage[a5paper, top=2cm, bottom=3cm, inner=2.2cm, outer=3cm]{geometry}
\usepackage[table]{xcolor}
\usepackage[italian]{babel}

% Colori
\definecolor{verdescuro}{HTML}{2d6a4f}
\definecolor{verdecharo}{HTML}{52b788}
\definecolor{grigioscuro}{HTML}{2b2d42}
\definecolor{giallo}{HTML}{d4a017}
\definecolor{rosso}{HTML}{c1121f}
\definecolor{blinfo}{HTML}{3a86ff}

\usepackage{tcolorbox}
\tcbuselibrary{skins,breakable}

\usepackage{booktabs}
\usepackage{tabularx}
\newcolumntype{L}{>{\raggedright\arraybackslash}X}

\usepackage{hyperref}
\hypersetup{
  colorlinks=true,
  linkcolor=verdescuro,
  urlcolor=verdescuro,
  pdftitle={Mangia Consapevole},
  pdfauthor={faustobe},
}

\usepackage{qrcode}
\usepackage{fancyhdr}
\usepackage{titlesec}
\usepackage{microtype}

% Riquadri colorati
\tcbset{
  boxbase/.style={
    breakable, enhanced, arc=4pt, boxrule=0.8pt,
    fonttitle=\bfseries\small,
    before upper={\parskip=4pt},
  }
}

\newcommand{\boxverde}[2]{%
  \begin{tcolorbox}[boxbase,
    colback=verdescuro!8!white, colframe=verdescuro, title={#1}]
  #2\end{tcolorbox}}

\newcommand{\boxgiallo}[2]{%
  \begin{tcolorbox}[boxbase,
    colback=giallo!12!white, colframe=giallo!80!black, title={#1}]
  #2\end{tcolorbox}}

\newcommand{\boxrosso}[2]{%
  \begin{tcolorbox}[boxbase,
    colback=rosso!8!white, colframe=rosso, title={#1}]
  #2\end{tcolorbox}}

\newcommand{\boxinfo}[2]{%
  \begin{tcolorbox}[boxbase,
    colback=blinfo!8!white, colframe=blinfo, title={#1}]
  #2\end{tcolorbox}}

% Stile pagina
\pagestyle{fancy}
\fancyhf{}
\fancyfoot[C]{\footnotesize\textcolor{grigioscuro}{\textit{Mangia Consapevole\ \textbar\ faustobe.it}}}
\fancyfoot[R]{\footnotesize\thepage}
\renewcommand{\headrulewidth}{0pt}
\renewcommand{\footrulewidth}{0.4pt}
\fancypagestyle{plain}{%
  \fancyhf{}%
  \fancyfoot[C]{\footnotesize\textcolor{grigioscuro}{\textit{Mangia Consapevole\ \textbar\ faustobe.it}}}%
  \fancyfoot[R]{\footnotesize\thepage}%
  \renewcommand{\headrulewidth}{0pt}%
  \renewcommand{\footrulewidth}{0.4pt}%
}

% Stile titoli capitolo
\titleformat{\chapter}[display]
  {\normalfont\LARGE\bfseries\color{verdescuro}}
  {\chaptertitlename\ \thechapter}{10pt}{\Huge}
\titlespacing*{\chapter}{0pt}{-20pt}{30pt}

% ============================================================
\begin{document}

% -------- COPERTINA --------
\begin{titlepage}
\pagecolor{verdescuro}
\color{white}
\vspace*{2.5cm}
\begin{center}
  {\fontsize{38}{44}\selectfont\bfseries Mangia Consapevole}\\[1.2cm]
  \textcolor{verdecharo}{\rule{6cm}{0.6pt}}\\[1.2cm]
  {\large\itshape Guida pratica agli additivi alimentari}\\[0.6cm]
  {\normalsize Con E-Codes Reader per Android}\\[3cm]
  {\normalsize\textcolor{verdecharo}{faustobe.it}}
\end{center}
\vfill
\begin{center}
  {\small Versione 1.0\ \textbullet\ 2026\ \textbullet\ CC BY-NC 4.0}
\end{center}
\end{titlepage}
\nopagecolor

\frontmatter
\tableofcontents
\mainmatter

% ============================================================
\chapter{Come usare l'app per fare la spesa consapevole}

Impara a riconoscere gli alimenti più lavorati e scegliere quelli
veramente semplici e naturali in pochi minuti, senza diventare
un esperto di additivi.

\section{Prima di scansionare, fai tre controlli rapidi}

Prima ancora di aprire l'app, osserva la confezione. Le etichette
frontali sono spesso pensate per sembrare più sani di quello che sono.

\begin{itemize}
  \item \textbf{Guarda il fronte della confezione.} Fai attenzione ai
    claim come ``senza zucchero ma con dolcificanti'', ``senza grassi ma
    con emulsionanti'', ``gusto extra'' o ``aromatizzato''. Spesso indicano
    che il prodotto è stato molto modificato.
  \item \textbf{Apri l'app e scansiona il prodotto.} Usa la fotocamera per
    scansionare il codice a barre o carica una foto dell'etichetta.
    L'app ti mostrerà subito la lista degli ingredienti e degli additivi.
  \item \textbf{Controlla tre punti chiave.} Quanti additivi compaiono
    (soprattutto i codici~E)? Quanto è lunga e complessa la lista
    ingredienti? Quanti zuccheri, sale e grassi sono dichiarati?
\end{itemize}

Se rispondi ``molti'' o ``lunga/complicata'' a più di una domanda, ti
trovi probabilmente davanti a un cibo molto lavorato.

\section{Tre regole per capire subito se un cibo è ``troppo lavorato''}

\textbf{Pochi ingredienti = meglio.} Se la lista è breve e contiene nomi
familiari (es.\ ``farina, acqua, sale, olio''), il prodotto è poco lavorato
e in genere più affidabile.

\textbf{Molti ingredienti ``impronunciabili''.} Termini come ``amidi modificati'',
``proteine isolate'', ``mono- e digliceridi degli acidi grassi'',
``esaltatori di sapidità'' o ``aromi complessi'' indicano chiaramente un
alimento molto processato.

\textbf{Tanti codici~E e additivi.} Non tutti gli~E sono pericolosi, ma
la loro presenza in massa --- coloranti, conservanti, stabilizzanti,
emulsionanti --- indica un grado elevato di trasformazione industriale.

\section{Come interpretare il risultato dell'app}

Ogni prodotto riceve una valutazione rapida: pensa a questi colori come
a un semaforo alimentare.

\boxverde{Poco lavorato --- via libera}{Pochi ingredienti, pochi o nessun
additivo~E. Zuccheri e sale moderati. Ottimo candidato per la spesa
quotidiana.}

\boxgiallo{Moderatamente lavorato --- con moderazione}{Alcuni~E, aromi
naturali o zuccheri extra. Va bene ogni tanto, non come base di ogni
pasto.}

\boxrosso{Ultra-processato --- limita o evita}{Molti additivi, zuccheri
extra, aromi complessi e grassi di origine industriale. Meglio tenerlo come
vera eccezione, specialmente per i bambini.}

\section{Come usare l'app in modo pratico in spesa}

\textbf{1.~Inizia dal reparto frigo e colazione.} Scansiona yogurt, pane,
latte, bevande. Cerca prodotti con pochi ingredienti, pochi~E, e zuccheri
semplici (frutta, latte, zucchero di canna) invece di sciroppi e
dolcificanti.

\textbf{2.~Poi spostati ai dolci e snack.} Se la lista è lunghissima e
piena di~E, l'app segnalerà rosso. In quei casi, scegli sempre
l'alternativa più semplice (frutta, pane integrale, noci).

\textbf{3.~Usa l'app per confrontare prodotti simili.} Cerca due prodotti
dello stesso tipo nell'app. Spesso quello con meno ingredienti è solo
leggermente meno appariscente sullo scaffale, ma molto più naturale.

\textbf{4.~Lascia i ``rossi'' come eccezione, non come regola.} Non serve
eliminare tutto ciò che è ultra-processato: basta ridurlo. Deciditi una
soglia --- ad esempio, massimo 1--2 prodotti ``rossi'' alla settimana ---
mentre il resto della spesa sia fatto di cibi semplici.

\section{Tabella riepilogativa}

{\small
\begin{tabularx}{\textwidth}{|L|L|}
\hline
\rowcolor{verdescuro}\textcolor{white}{\textbf{Situazione nell'app}}
  & \textcolor{white}{\textbf{Cosa fare in pratica}} \\
\hline
Pochi ingredienti, pochi/no~E
  & Base quotidiana: yogurt semplice, pane integrale, passata di pomodoro \\
\hline
Alcuni~E, zuccheri moderati
  & Ogni tanto, non come alimento principale \\
\hline
Tanti~E, zuccheri extra, aromi complessi
  & Limita o evita: usalo solo come eccezione \\
\hline
\end{tabularx}
}

\section{Perché imparare a leggere gli additivi fa davvero la differenza}

Riconoscere i cibi più lavorati non è ``fobia degli additivi'': è
\textbf{attenzione alla qualità di quello che mangiamo}. Molti studi
suggeriscono che ridurre i cibi ultra-processati aiuta sul lungo periodo
il controllo del peso, la salute cardiovascolare e la stabilità del
metabolismo.

L'app ti dà uno strumento concreto: non serve ricordare ogni codice~E,
basta imparare a riconoscere i segnali di troppa trasformazione.

% ============================================================
\chapter{Come leggere l'etichetta nutrizionale}

Ingredienti, calorie, sale nascosto: impara a decifrare un'etichetta in
meno di un minuto e smetti di comprare prodotti che sembrano sani ma non
lo sono.

\section{La lista degli ingredienti: dove guardare per prima cosa}

Per legge, gli ingredienti sono elencati in \textbf{ordine decrescente di
peso}: ciò che compare prima è presente in quantità maggiore. È il punto
di partenza più importante dell'etichetta.

\begin{itemize}
  \item \textbf{Il primo ingrediente conta di più.} Se il primo
    ingrediente è ``zucchero'', ``sciroppo di glucosio'' o ``olio di palma'',
    il prodotto è costruito attorno a quell'elemento.
  \item \textbf{Lunghezza della lista = grado di lavorazione.} Pane fatto
    in casa: farina, acqua, sale, lievito --- 4 ingredienti. Pane
    industriale da scaffale: può averne 15 o più.
  \item \textbf{Cerca i nomi difficili da pronunciare.} Termini come
    ``amido modificato'', ``mono- e digliceridi degli acidi grassi'',
    ``carboxymethyl cellulose'' o ``idrogenato'' indicano ingredienti di
    origine industriale.
\end{itemize}

\boxinfo{Regola pratica}{Se non riusciresti a trovare tutti gli ingredienti
in un negozio di alimentari normale, probabilmente il prodotto è molto
lavorato.}

\section{La tabella nutrizionale: cosa guardare (e cosa ignorare)}

\begin{description}
  \item[Zuccheri] Meno di 5\,g per 100\,g è basso; tra 5 e 22,5\,g è
    medio; sopra 22,5\,g è alto. Per le bevande il limite ``basso'' è
    2,5\,g/100\,ml.
  \item[Sale] Meno di 0,3\,g/100\,g è basso; sopra 1,5\,g/100\,g è
    alto. Il sodio va moltiplicato per 2,5 per ottenere il sale.
  \item[Grassi saturi] Meno di 1,5\,g/100\,g è accettabile; sopra
    5\,g/100\,g è considerato alto.
\end{description}

\section{Porzioni vs 100\,g: il trucco che confonde}

I valori nutrizionali vengono spesso indicati sia per 100\,g che ``per
porzione''. Le aziende tendono a mettere in evidenza i valori per porzione
--- ma le porzioni sono spesso sottostimate.

\boxgiallo{Esempio}{Se la confezione dichiara 10\,g di zuccheri per
porzione (30\,g) e tu ne mangi una ciotola da 60\,g, stai assumendo 20\,g
di zuccheri --- quasi la metà della soglia giornaliera raccomandata.}

\section{Il sale nascosto: dove si nasconde davvero}

Il sale non si trova solo nei prodotti salati.

\begin{itemize}
  \item \textbf{Pane e prodotti da forno:} due fette di pane industriale
    possono contenere 0,5--0,8\,g di sale.
  \item \textbf{Cereali da colazione:} anche quelli ``sani'' e ``integrali''
    possono contenere 0,8--1,2\,g di sale per 100\,g.
  \item \textbf{Salse, sughi e condimenti:} possono contenere tra 1 e
    3\,g di sale per 100\,g.
\end{itemize}

\boxinfo{Nota bene}{In etichetta potresti trovare ``sodio'' invece di
``sale''. Per convertire: \textbf{sale = sodio \texttimes\ 2,5}.}

\section{I claim nutrizionali: cosa significano davvero}

I claim come ``senza grassi'', ``light'', ``ricco di fibre'' o ``0\% zuccheri''
sono regolamentati dalla legge, ma possono trarre in inganno.

\begin{itemize}
  \item \textbf{``Senza zuccheri aggiunti''} non significa senza zuccheri:
    può contenerne di naturali (fruttosio, lattosio) o dolcificanti
    artificiali.
  \item \textbf{``Light'' o ``ridotto''} significa ridotto del 30\%
    rispetto al prodotto di riferimento. Spesso il grasso tolto è
    sostituito con zuccheri o addensanti.
  \item \textbf{``Ricco di fibre''} richiede almeno 6\,g/100\,g ma non
    dice nulla sugli altri ingredienti: un biscotto ``ricco di fibre''
    può contenere molto zucchero.
\end{itemize}

\section{Checklist rapida da usare al supermercato}

\begin{enumerate}
  \item Guarda il primo ingrediente: è un alimento vero?
  \item Conta gli ingredienti: meno di 5 è ottimo; oltre 10--12 controlla
    bene.
  \item Controlla i codici~E: pochi o nessuno è un buon segnale.
  \item Guarda sale e zuccheri per 100\,g: sale oltre 1,5\,g o zuccheri
    oltre 22,5\,g sono valori alti.
  \item Ignora i claim promozionali sul fronte: fida sempre della lista
    ingredienti.
\end{enumerate}

% ============================================================
\chapter{Additivi vs ultra-processati: qual è la differenza?}

Avere molti codici~E non è la stessa cosa che essere ultra-processato.
Capire la differenza ti aiuta a fare valutazioni più complete e precise.

\section{Cosa sono gli additivi alimentari}

Gli additivi alimentari sono sostanze aggiunte intenzionalmente agli
alimenti per svolgere una funzione tecnologica: conservare, colorare,
addensare, emulsionare, stabilizzare o migliorare il sapore. In Europa,
ogni additivo approvato riceve un codice \textbf{E} e deve superare una
valutazione di sicurezza dell'EFSA.

\begin{itemize}
  \item \textbf{Non tutti gli additivi sono sintetici.} L'acido citrico
    (E330), la curcumina (E100), la lecitina di soia (E322) e l'acido
    ascorbico (E300) sono additivi di origine naturale.
  \item \textbf{Non tutti gli additivi sono pericolosi.} La stragrande
    maggioranza degli~E approvati è considerata sicura alle dosi
    consentite.
  \item \textbf{Alcuni additivi meritano attenzione.} Coloranti azoici
    (E102, E110, E122, E124, E129), nitrati (E249--E252), dolcificanti
    intensivi (E950--E962): queste categorie hanno evidenze di rischio
    più consistenti.
\end{itemize}

\section{Cosa sono i cibi ultra-processati}

Il termine ``ultra-processato'' viene dalla \textbf{classificazione NOVA},
sviluppata dall'Università di San Paolo. Non riguarda i nutrienti, ma il
\emph{processo industriale} che ha prodotto l'alimento.

Un cibo è ultra-processato (NOVA 4) quando contiene ingredienti che non
si trovano in una cucina domestica: proteine isolate, amidi modificati,
sciroppi di glucosio-fruttosio, aromi ``natura-identici'', emulsionanti,
coloranti, stabilizzanti --- spesso in combinazioni complesse.

\section{Quando coincidono --- e quando no}

\boxverde{Additivi ma NON ultra-processato}{Formaggio con conservante E252:
lavorato ma non NOVA~4. Vino con solfiti (E220): additivo naturale,
prodotto tradizionale. Olio con vitamina~E (E306): minima lavorazione.}

\boxrosso{Ultra-processato con POCHI additivi~E}{Pane da toast con amido
modificato e aromi: non ha tanti~E ma è NOVA~4. Barrette proteiche con
proteine isolate: pochi codici, ma fortemente trasformate.}

\boxgiallo{Ultra-processato E con tanti additivi (il caso peggiore)}{Merendine:
amidi modificati + coloranti + emulsionanti + aromi + dolcificanti. Salse
industriali: conservanti + stabilizzanti + esaltatori di sapidità.}

\section{Il doppio problema: perché entrambi contano}

\begin{itemize}
  \item \textbf{Rischi legati agli additivi specifici:} i coloranti azoici
    e l'iperattività nei bambini (EFSA 2007); i nitrati nella carne
    lavorata e le nitrosammine cancerogene; i dolcificanti intensivi e la
    possibile alterazione del microbioma intestinale.
  \item \textbf{Rischi legati all'ultra-lavorazione in sé:} studi su larga
    scala (NutriNet-Santé, UK~Biobank) associano il consumo elevato di
    NOVA~4 a obesità, diabete tipo~2, malattie cardiovascolari e alcuni
    tumori --- indipendentemente dalla qualità nutrizionale.
  \item \textbf{L'effetto combinazione:} gli studi sugli additivi testano
    quasi sempre una sostanza alla volta. Nella vita reale mangiamo
    combinazioni di 10--20 additivi diversi nello stesso pasto.
\end{itemize}

\section{Come riconoscerli in pratica}

\begin{enumerate}
  \item Guarda se ci sono ingredienti ``da laboratorio'': amidi modificati,
    proteine isolate, sciroppi di glucosio-fruttosio, aromi
    ``natura-identici''.
  \item Conta i codici~E e identifica le categorie critiche: un~E
    antiossidante naturale (E306) è molto diverso da un colorante azoico
    (E102).
  \item Considera il contesto del pasto: un pasto equilibrato con verdure
    e legumi ``assorbe'' meglio qualche prodotto più lavorato.
  \item Usa l'app come supporto, non come giudice assoluto.
\end{enumerate}

\boxinfo{Riassunto in una frase}{Un prodotto con un~E naturale non è un
problema. Un prodotto con 15 ingredienti industriali senza nessun~E è
comunque un ultra-processato. Guarda entrambe le cose.}

% ============================================================
\chapter{Classificazione NOVA: cos'è e come usarla}

Il sistema NOVA classifica gli alimenti in base al grado di trasformazione
industriale --- non alle calorie o ai nutrienti. Capire i quattro gruppi
ti aiuta a fare scelte migliori ogni giorno.

\section{Cos'è la classificazione NOVA}

NOVA è un sistema sviluppato dal professor Carlos Monteiro all'Università
di San Paolo (Brasile) e oggi adottato dall'OMS, dalla FAO e da numerosi
studi epidemiologici globali.

A differenza dei sistemi basati sui nutrienti (come il Nutri-Score),
\textbf{NOVA classifica in base al grado di trasformazione industriale}:
non ``quante calorie ha'', ma ``quanta industria è entrata nella sua
produzione''.

\section{Gruppo 1 --- Cibi non lavorati o minimamente lavorati}

Sono alimenti di origine vegetale o animale che non hanno subito nessuna
trasformazione industriale significativa, o solo processi fisici minimi
(essiccazione, macinazione, raffreddamento, pastorizzazione).

\textbf{Esempi:} frutta e verdura fresca o congelata, carne e pesce
freschi, uova, legumi secchi, cereali in chicco (riso, avena, farro),
latte fresco pastorizzato, yogurt naturale senza additivi.

I cibi del Gruppo~1 dovrebbero costituire la base della dieta.

\section{Gruppo 2 --- Ingredienti culinari lavorati}

Sono sostanze estratte da alimenti del Gruppo~1 o dalla natura, usate in
cucina per cucinare e condire. Non si mangiano da soli.

\textbf{Esempi:} olio d'oliva, burro, aceto, sale, zucchero, miele,
farina, amido di mais puro. Usali con moderazione.

\section{Gruppo 3 --- Cibi lavorati}

Sono prodotti ottenuti aggiungendo ingredienti del Gruppo~2 ad alimenti
del Gruppo~1 (conservazione sotto sale, affumicatura, fermentazione,
stagionatura). Riconoscibili perché contengono solo 2--3 ingredienti.

\textbf{Esempi:} verdure in conserva, formaggi tradizionali (parmigiano,
pecorino), prosciutto crudo di qualità, pane artigianale, vino.

\section{Gruppo 4 --- Ultra-processati: il gruppo da limitare}

Formulazioni industriali che non contengono alimenti interi o li usano
in minima parte. Il segnale chiave: contengono ingredienti che \textbf{non
si trovano in una cucina domestica}.

\textbf{Ingredienti tipici:} proteine idrolizzate, amidi modificati,
sciroppo di glucosio-fruttosio, oli idrogenati, aromi ``natura-identici'',
coloranti artificiali, dolcificanti intensivi, emulsionanti (E471, E472).

\boxrosso{Dati chiave}{In Italia circa il 14--18\% delle calorie
giornaliere proviene da ultra-processati. In UK e USA la percentuale
supera il 50--60\%. Già un consumo superiore al 20\% delle calorie da
NOVA~4 è associato a rischi per la salute.}

\section{NOVA nella spesa pratica}

{\small
\begin{tabularx}{\textwidth}{|L|L|L|L|}
\hline
\rowcolor{verdescuro}
\textcolor{white}{\textbf{Gruppo}}
  & \textcolor{white}{\textbf{Definizione}}
  & \textcolor{white}{\textbf{Esempi}}
  & \textcolor{white}{\textbf{In dieta}} \\
\hline
NOVA~1 & Non/min.\ lavorati & Frutta, verdura, uova, carne fresca & Base \\
\hline
NOVA~2 & Ingredienti culinari & Olio, sale, zucchero, farina & Con moderazione \\
\hline
NOVA~3 & Lavorati tradizionali & Conserve, formaggi, pane artigianale & Accettabile \\
\hline
NOVA~4 & Ultra-processati & Merendine, wurstel, cereali dolci & Limitare \\
\hline
\end{tabularx}
}

% ============================================================
\chapter{Ingredienti da monitorare sulle etichette}

Zuccheri nascosti, grassi pericolosi, additivi controversi e sodio in
eccesso: impara a riconoscerli a colpo d'occhio.

\section{Zuccheri: i 10 nomi nascosti sulle etichette}

Lo zucchero non si chiama sempre ``zucchero''. Le aziende usano decine di
varianti per distribuirlo in più voci della lista ingredienti.

\begin{itemize}
  \item \textbf{Sciroppo di glucosio-fruttosio} (il peggiore)
  \item \textbf{Maltodestrine} (amido quasi interamente scomposto)
  \item \textbf{Destrosio}, \textbf{Fruttosio}
  \item \textbf{Succo di canna concentrato}, \textbf{Sciroppo di mais}
  \item \textbf{Sciroppo di riso}, \textbf{Miele disidratato}
  \item \textbf{Zucchero di cocco}, \textbf{Nettare di agave}
\end{itemize}

L'OMS raccomanda di mantenere gli zuccheri liberi sotto il \textbf{10\%
delle calorie giornaliere} --- circa 50\,g al giorno per un adulto da
2\,000\,kcal. Preferibilmente sotto il 5\% (25\,g).

\section{Grassi da evitare: idrogenati e trans}

Gli acidi grassi trans, prodotti dall'idrogenazione industriale degli oli
vegetali, sono associati a un aumentato rischio cardiovascolare.

\boxrosso{Grassi idrogenati e parzialmente idrogenati}{Cerca in etichetta:
``olio vegetale parzialmente idrogenato'', ``grasso vegetale idrogenato'',
``shortening''. In Europa sono vietati sopra 2\,g/100\,g di grasso dal
2021. Evitali comunque.}

\boxgiallo{Olio di palma: non è trans, ma\ldots}{Durante la raffinazione
ad alte temperature si formano composti (3-MCPD, glicidolo) potenzialmente
cancerogeni. Non è vietato, ma meglio limitarlo.}

\section{Additivi controversi: i codici~E da conoscere}

\begin{itemize}
  \item \textbf{Coloranti azoici --- massima attenzione con i bambini:}
    E102 (Tartrazina), E104, E110, E122, E124, E129. Studio EFSA 2007:
    correlazione con iperattività nei bambini. In Europa obbligo di avviso:
    ``può influire sull'attività e l'attenzione dei bambini''.
  \item \textbf{Nitrati e nitriti --- carne lavorata:} E249--E252. In
    presenza di proteine animali formano nitrosammine, classificate come
    probabili cancerogeni (IARC gruppo~2A).
  \item \textbf{Dolcificanti intensivi:} E950 (Acesulfame K), E951
    (Aspartame), E952 (Ciclamato), E954 (Saccarina), E955 (Sucralosio).
    L'OMS (2023) ha espresso cautela sull'uso cronico.
\end{itemize}

\section{Sodio nascosto: oltre il sale da cucina}

Il sodio in eccesso alza la pressione arteriosa. Circa il 75\% del sodio
nella dieta occidentale viene già nel prodotto confezionato.

\textbf{Fonti principali:} pane, formaggi, salumi, piatti pronti, sughi
in barattolo, dado e brodi, salse tipo soia, ketchup, cereali da colazione.

Anche additivi come il glutammato monosodico (E621) e il benzoato di
sodio (E211) contribuiscono al sodio totale.

\section{La regola del 5: il filtro più veloce}

\begin{enumerate}
  \item Più di 5 ingredienti? Inizia a leggere con attenzione.
  \item Zuccheri sopra 5\,g/100\,g? Controlla da dove vengono.
  \item Sale sopra 0,5\,g/100\,g in un prodotto dolce? Segnale di
    ultra-lavorazione.
  \item Più di 5 additivi~E? Il prodotto è probabilmente ultra-processato.
  \item Non riconosci 5 o più ingredienti? Rimettilo sullo scaffale.
\end{enumerate}

% ============================================================
\chapter{Alimenti consigliati}

Sostituire non eliminare: le alternative semplici, accessibili e spesso
più economiche agli ultra-processati più comuni.

\section{Il principio base: sostituire, non eliminare}

Non si tratta di fare diete rigide. Si tratta di conoscere le alternative
--- e di sceglierle quando sono disponibili e pratiche quanto il prodotto
che sostituiscono.

\textbf{La regola 80/20 funziona:} se l'80\% di quello che mangi è
Gruppo~1--2 (NOVA), il 20\% di prodotti più lavorati non fa la differenza.
La salute alimentare si misura sulla media settimanale, non sul singolo
pasto.

\section{Cereali e carboidrati: le scelte migliori}

\textbf{Cereali consigliati:} avena in fiocchi (senza ingredienti
aggiunti), riso integrale, farro in chicco, orzo perlato, grano saraceno,
pasta integrale con sola semola integrale, pane di sola farina integrale
(massimo 4--5 ingredienti).

\textbf{Da limitare:} cereali da colazione con zucchero, pane da cassetta
con amidi modificati, cracker con olio di palma, riso e pasta pre-cotti
in busta con additivi.

\boxinfo{Tip rapido per i cereali da colazione}{Se i zuccheri sono sopra
10\,g/100\,g, stai mangiando un prodotto da dessert --- non una colazione
sana.}

\section{Proteine: le fonti migliori}

\begin{itemize}
  \item \textbf{Pesce e frutti di mare:} pesce fresco o surgelato senza
    impanatura, sardine e sgombri in olio d'oliva, tonno al naturale,
    acciughe.
  \item \textbf{Uova e latticini:} uova intere, yogurt greco naturale,
    ricotta, fiocchi di latte, formaggi stagionati con moderazione.
  \item \textbf{Legumi:} lenticchie, ceci, fagioli, piselli. In scatola
    al naturale sono già cotti e si aggiungono in 30 secondi a qualsiasi
    pasto.
\end{itemize}

\section{Grassi buoni: i più semplici da aggiungere}

\textbf{Fonti consigliate:} olio d'oliva extravergine, olio di lino crudo,
noci, mandorle, nocciole, semi di lino, semi di chia, avocado, pesce
azzurro, uova intere.

\textbf{Con moderazione:} burro, olio di cocco, latte intero.

\textbf{Da evitare:} margarine con grassi idrogenati, shortening vegetale,
panna vegetale con emulsionanti.

\section{Lista della spesa tipo}

{\small
\begin{tabularx}{\textwidth}{|L|L|}
\hline
\rowcolor{verdescuro}\textcolor{white}{\textbf{Categoria}}
  & \textcolor{white}{\textbf{Prodotti consigliati}} \\
\hline
Cereali & Avena, riso integrale, farro, pasta integrale, pane artigianale \\
\hline
Legumi & Lenticchie, ceci, fagioli (secchi o in scatola al naturale) \\
\hline
Proteine animali & Uova, sardine, tonno al naturale, yogurt greco, ricotta \\
\hline
Verdure & Di stagione, fresche o surgelate senza additivi \\
\hline
Grassi & Olio EVO, noci, mandorle, semi di lino \\
\hline
Frutta & Intera di stagione, frutti di bosco surgelati senza zucchero \\
\hline
Condimenti & Aceto di mele, salsa di soia a bassa lavorazione, spezie \\
\hline
\end{tabularx}
}

% ============================================================
\chapter{Perché gli alimenti ultra-processati fanno male}

Le evidenze scientifiche sui rischi per la salute legati al consumo
abituale di ultra-processati: obesità, infiammazione, microbiota,
malattie croniche.

\section{Obesità e alterazione del metabolismo}

Gli alimenti ultra-processati non fanno ingrassare semplicemente perché
sono calorici. Il meccanismo è più insidioso: sono progettati per essere
\textbf{iper-appetibili} e interferiscono attivamente con i segnali
naturali di sazietà.

\boxinfo{Studio di Kevin Hall (NIH, 2019)}{In un trial clinico controllato,
i partecipanti esposti a una dieta a base di ultra-processati consumavano
in media \textbf{500 calorie in più al giorno} rispetto a chi mangiava
cibi minimamente lavorati --- a parità di libertà di scelta. In due
settimane, il gruppo UPF aveva guadagnato quasi un chilo.}

\begin{itemize}
  \item \textbf{Alta densità energetica:} gli UPF contengono in media
    2,15\,kcal per grammo --- quasi il doppio rispetto ai cibi freschi.
  \item \textbf{Stimolazione eccessiva dei centri del piacere:} la
    combinazione di zucchero, sale, grassi e aromi attiva il sistema
    dopaminergico. Si mangia anche in assenza di fame reale.
  \item \textbf{Riduzione del peptide YY:} gli UPF abbassano i livelli
    dell'ormone intestinale che segnala la sazietà al cervello.
\end{itemize}

\boxrosso{+41\% di rischio di obesità addominale}{Chi segue una dieta ad
alto contenuto di UPF presenta un rischio superiore del 41\% di sviluppare
obesità rispetto a chi ne consuma pochi.}

\section{Infiammazioni e microbiota intestinale}

Il consumo frequente di ultra-processati è associato a un aumento delle
risposte immunitarie infiammatorie. Emulsionanti come la
carbossimetilcellulosa (E466) e il polisorbato~80 (E433) alterano il
muco protettivo intestinale anche a dosi consentite.

Il processo di ultra-trasformazione priva gli alimenti di fibre e composti
vegetali bioattivi, favorendo la \textbf{disbiosi} --- uno squilibrio del
microbiota intestinale associato a insulino-resistenza e infiammazione
cronica.

\section{Un rischio sistemico: più meccanismi insieme}

{\small
\begin{tabularx}{\textwidth}{|L|L|L|}
\hline
\rowcolor{verdescuro}
\textcolor{white}{\textbf{Meccanismo}}
  & \textcolor{white}{\textbf{Effetto}}
  & \textcolor{white}{\textbf{Conseguenza}} \\
\hline
Iper-appetibilità & Consumo eccessivo & Obesità, diabete tipo~2 \\
\hline
Carenza di fibre & Disbiosi intestinale & Infiammazione cronica, IBD \\
\hline
Additivi (emulsionanti) & Permeabilità intestinale & Malattie autoimmuni \\
\hline
Zuccheri e dolcificanti & Insulino-resistenza & Diabete tipo~2 \\
\hline
Sostituzione del cibo fresco & Carenza micronutrienti & Stress ossidativo \\
\hline
\end{tabularx}
}

\boxinfo{I grandi studi epidemiologici}{NutriNet-Santé (Francia, oltre
100.000 partecipanti) e UK~Biobank (UK, oltre 500.000 partecipanti)
associano il consumo elevato di UPF a maggiori rischi di cancro, malattie
cardiovascolari, depressione e mortalità per tutte le cause ---
indipendentemente dalla qualità nutrizionale dei singoli prodotti.}

\section{Cosa fare in pratica}

\begin{enumerate}
  \item \textbf{Individua i tuoi UPF abituali:} cereali dolci, snack
    confezionati, pane da cassetta, piatti pronti, bevande aromatizzate.
  \item \textbf{Sostituisci uno alla volta:} avena al posto dei cereali
    dolci, yogurt greco al posto di quello aromatizzato.
  \item \textbf{Aumenta le fibre --- sono la priorità:} legumi, verdure,
    cereali integrali, frutta intera.
  \item \textbf{Usa l'app per identificare gli UPF nascosti:} molti
    prodotti che sembrano sani sono in realtà NOVA~4.
  \item \textbf{Non puntare alla perfezione:} anche un miglioramento
    del 20--30\% nella qualità della dieta produce effetti misurabili
    sulla salute nel medio termine.
\end{enumerate}

% ============================================================
\chapter{Sicurezza alimentare per bambini e anziani}

Stessa dose, effetti diversi: perché le categorie più vulnerabili
richiedono attenzione particolare agli additivi e al sodio.

\section{Perché bambini e anziani sono più vulnerabili}

La stessa dose di additivo che un adulto sano metabolizza senza problemi
può avere effetti diversi su un bambino piccolo o su una persona anziana.

\begin{itemize}
  \item \textbf{Bambini: peso corporeo basso, sistemi non maturi.} La dose
    giornaliera accettabile (DGA) degli additivi è calcolata in mg per kg
    di peso corporeo. Un bambino da 15\,kg che mangia una merenda da 30\,g
    con coloranti riceve una dose relativa 4--5 volte superiore a quella
    di un adulto da 70\,kg.
  \item \textbf{Anziani: ridotta funzionalità renale ed epatica.} Con
    l'età diminuisce la capacità di eliminare alcune sostanze. Il sodio
    in eccesso impatta di più su chi soffre di ipertensione o
    insufficienza renale.
  \item \textbf{Effetti cumulativi:} non è il singolo gelato con coloranti
    una volta al mese il problema. Se i bambini mangiano quotidianamente
    prodotti con E102, E110, E122 a colazione, a merenda, a pranzo ---
    l'effetto cumulativo diventa rilevante.
\end{itemize}

\section{Additivi critici per i bambini}

\boxrosso{Coloranti azoici --- da evitare con i bambini}{E102 (Tartrazina),
E104, E110, E122, E124, E129. In Europa obbligo di etichetta: \emph{``può
influire sull'attività e sull'attenzione dei bambini''}. Si trovano in
caramelle colorate, bibite arancioni/rosse, gelatine, dolci industriali.}

\boxgiallo{Dolcificanti intensivi --- non approvati per bambini piccoli}{E951
(Aspartame), E950 (Acesulfame K), E955 (Sucralosio), E954 (Saccarina).
Non autorizzati negli alimenti per lattanti (Reg.\ UE 1333/2008). Si
trovano in yogurt ``0\% zucchero'', bibite ``light'', caramelle senza
zucchero.}

\section{Cosa limitare negli anziani}

\begin{itemize}
  \item \textbf{Sodio: il rischio principale.} I piatti pronti, i salumi,
    i formaggi stagionati e i brodi industriali sono i principali veicoli.
    Obiettivo: meno di 5\,g di sale al giorno.
  \item \textbf{Fosfati: interferiscono con calcio e ossa.} E338--E341,
    E450--E452 sono comuni nei prodotti a base di carne lavorata e
    formaggi fusi. Un consumo elevato interferisce con l'assorbimento del
    calcio.
  \item \textbf{Interazioni con farmaci:} la vitamina~K nei vegetali con
    i warfarinici, il pompelmo con le statine. Non si tratta di evitare
    questi alimenti, ma di mantenere la dieta costante.
\end{itemize}

\section{Lista sicura per la spesa}

{\small
\begin{tabularx}{\textwidth}{|L|L|L|}
\hline
\rowcolor{verdescuro}
\textcolor{white}{\textbf{Additivo/Ingrediente}}
  & \textcolor{white}{\textbf{Bambini}}
  & \textcolor{white}{\textbf{Anziani}} \\
\hline
Coloranti azoici (E102, E110, E122, E124, E129) & Evitare & Limitare \\
\hline
Dolcificanti intensivi (aspartame, sucralosio\ldots) & Non per bambini piccoli & Con moderazione \\
\hline
Sodio elevato (\textgreater1,5\,g sale/100\,g) & Limitare & Evitare \\
\hline
Fosfati (E338--E452) & Moderazione & Limitare \\
\hline
Caffeina (energy drink, soft drink) & Evitare \textless12 anni & Attenzione \\
\hline
Nitrati/nitriti (E249--E252) in salumi & Limitare & Limitare \\
\hline
\end{tabularx}
}

% ============================================================
% Pagina finale con QR code
\cleardoublepage
\thispagestyle{empty}
\begin{center}
\vspace*{3cm}
{\large\bfseries\color{verdescuro} Scarica E-Codes Reader}\\[0.8cm]
{\normalsize Scansiona il QR code per trovare l'app\\
sul Google Play Store.}\\[1.5cm]
\qrcode[height=3.5cm]{https://play.google.com/store/apps/details?id=com.faustobe.ecodes}\\[1.5cm]
\textcolor{grigioscuro}{\footnotesize\textit{faustobe.it\ \textbullet\ Versione 1.0\ \textbullet\ 2026\ \textbullet\ CC BY-NC 4.0}}
\end{center}

\end{document}
